% ============================================================================
% 4.2 adder.v
% ============================================================================
\subsection{设计目标}

把 64 位输入（8 个字节）的每个字节提取出来，求和，输出一个 8 位结果。
这是一个纯组合逻辑（不需要时钟）+ 一个寄存器输出的简单示例。

\subsection{源码与讲解}

\begin{lstlisting}[style=verilog,caption=Chip/src/adder.v]
module adder
#(
    parameter CLK_FRE = 50,     // 时钟频率(MHz)
    parameter BAUD_RATE = 115200 // 波特率（这里未使用，但保持与其他模块一致）
)
(
    input clk,                   // 时钟输入
    input rst_n,                 // 异步复位，低电平有效
    input [63:0] in_register,    // 64 位输入（8 个字节）
    output reg [7:0] out_register // 8 位输出（求和结果）
);

always @(posedge clk or negedge rst_n)
begin
    if (!rst_n)
    begin
        out_register <= 8'd0;
    end
    else
    begin
        out_register
            <= (in_register >> (8*0) & 8'b11111111)
             + (in_register >> (8*1) & 8'b11111111)
             + (in_register >> (8*2) & 8'b11111111)
             + (in_register >> (8*3) & 8'b11111111)
             + (in_register >> (8*4) & 8'b11111111)
             + (in_register >> (8*5) & 8'b11111111)
             + (in_register >> (8*6) & 8'b11111111)
             + (in_register >> (8*7) & 8'b11111111);
    end
end
endmodule
\end{lstlisting}

\paragraph{位提取与掩码}
\begin{itemize}
    \item \texttt{in\_register >> (8*0)} 不移位；\texttt{\& 8'hFF} 只保留最低字节
    \item \texttt{in\_register >> (8*1)} 右移 8 位，保留第二字节
    \item \ldots 直到 \texttt{8*7} = 第 8 个字节（最高位）
\end{itemize}

每个字节都是 8 位，8 个这样的数相加，最大可能值是
\( 8 \times 255 = 2040 \)，需要至少 11 位存。但我们的
\texttt{out\_register} 只有 8 位——结果会\textbf{溢出并截断}，
只保留低 8 位。这是故意的——它演示了"硬件里做的运算宽度是
你事先规定好的，超出部分就丢掉"的特点。

\paragraph{综合后会是什么样的硬件}
综合器会把这个设计变成一个由 7 个 8 位加法器串起来的电路
（加法树），最后一级的和被一个 8 位寄存器锁住。每个时钟沿，
寄存器都会重新计算一次——如果输入不变，输出也不变。

如果你想保留完整的 11 位结果，只需把 \texttt{out\_register}
改成 \texttt{[10:0]}（11 位）。其余代码无需改动。

\subsection{一个更简洁的写法}

其实 Verilog 支持用 \texttt{for} 循环写重复逻辑：

\begin{lstlisting}[style=verilog,caption=用 generate 写的简化版本]
// generate 块在综合时被展开成具体的门电路
generate
    for (i = 0; i < 8; i = i + 1)
    begin : gen_bytes
        // 编译时展开：8 个独立的字节提取逻辑
    end
endgenerate
\end{lstlisting}

但对我们的目标而言——理解硬件在做什么——
直接写 8 次反而更清晰，因为你能看到每个字节在做同一件事。
